\chapter{Literature Review}
\label{ch:sota}

The field of Explainable Artificial Intelligence (XAI) has expanded rapidly in response to the opacity of modern Deep Learning (DL) architectures. While early research focused primarily on generating visually appealing explanations, recent scholarship has shifted towards rigorous quantification of explanation quality. This chapter provides a taxonomy of Feature Attribution Explanation (FAE) methods, analyzes the "Crisis of Evaluation" precipitated by the lack of ground truth, and systematizes the mathematical frameworks used to measure Faithfulness, Robustness, and Complexity. Finally, it reviews the theoretical foundations of XAI Ensembling as a mechanism to mitigate the limitations of individual estimators.

\section{Feature Attribution Explanation (FAE) Methods}

Feature attribution methods, often referred to as saliency maps in computer vision, aim to assign a scalar relevance score $R_i$ to each input feature $x_i$ (e.g., a pixel), representing its contribution to the model's output $f(x)$ for a specific target class $c$. Formally, an explanation method $\Phi$ takes a model $f$ and an input $x$ to produce a map $E = \Phi(f, x)$, where $E \in \mathbb{R}^{d}$ and $d$ is the dimensionality of the input.

Current literature broadly categorizes these methods based on their access to the model's internal parameters: Gradient-based (White-box) and Perturbation-based (Black-box) approaches \cite{hedstrom2023quantus}.


\subsection{Gradient-based Methods (White-box)}
Gradient-based methods exploit the differentiable nature of neural networks. They require access to the model's weights and gradients, operating on the intuition that the gradient of the output with respect to the input represents the model's sensitivity to feature changes. The methods below form a progression: from raw gradients, to path-integrated gradients that restore axiomatic guarantees, to feature-space gradients that trade pixel resolution for class-discriminative localization.

\subsubsection{Saliency Maps}
The foundational approach, proposed by Simonyan et al., computes the gradient of the class score $S_c$ with respect to the input image $I$ \cite{simonyan2014deep}. The relevance map is defined as:
\begin{equation}
    E_{Vanilla} = \left| \frac{\partial S_c}{\partial I} \right|
\end{equation}
While computationally efficient, vanilla gradients often suffer from "shattered gradients," resulting in noisy visualizations that lack class-discriminativeness. Furthermore, they are susceptible to the "saturation problem," where a feature may strongly contribute to a prediction, but its gradient is zero because the activation is saturated (e.g., in a ReLU function) \cite{sundararajan2017axiomatic}.

\subsubsection{Integrated Gradients (IG)}
To address gradient saturation, Sundararajan et al. introduced Integrated Gradients, which fundamentally relies on the \textit{Axiom of Completeness}. IG accumulates gradients along a linear path from a baseline input $x'$ (usually a black image) to the actual input $x$ \cite{sundararajan2017axiomatic}. The attribution for the $i$-th feature is given by:
\begin{equation}
    IG_i(x) = (x_i - x'_i) \times \int_{\alpha=0}^{1} \frac{\partial f(x' + \alpha(x - x'))}{\partial x_i} d\alpha
\end{equation}
This method ensures that the sum of attributions equals the difference between the model output at the input and the baseline, providing a theoretical guarantee missing in vanilla gradients.

\subsubsection{Grad-CAM}
Gradient-weighted Class Activation Mapping (Grad-CAM) moves away from pixel-space gradients to feature-space gradients. It utilizes the gradients flowing into the final convolutional layer of a CNN to produce a coarse localization map \cite{selvaraju2017grad}. Global Average Pooling is applied to the gradients to obtain neuron importance weights $\alpha_k^c$:
\begin{equation}
    \alpha_k^c = \frac{1}{Z} \sum_i \sum_j \frac{\partial y^c}{\partial A_{ij}^k}
\end{equation}
where $A^k$ is the $k$-th feature map. The final map is a ReLU-activated linear combination of these maps. Grad-CAM is highly effective for localization but lacks fine-grained pixel detail, often necessitating combination with Guided Backpropagation (Guided Grad-CAM).

\subsection{Perturbation-based Methods (Black-box)}
Perturbation-based methods operate under a model-agnostic paradigm. They estimate feature importance by systematically modifying (perturbing) the input and observing the resulting changes in the model's output prediction. Because they treat the model as a black box, they apply uniformly across architectures, but their reliance on sampling makes them both computationally expensive and, as noted below, susceptible to estimation noise. The thesis surveys LIME and SHAP here for completeness, but---as Chapter~\ref{ch:methodology} details---the final evaluation panel retains only the gradient-based methods together with Occlusion, since the Robustness metrics re-invoke the explanation function for every perturbation sample, multiplying the already-substantial surrogate-fitting cost of LIME and KernelSHAP beyond what the scale of the full study permits.

\subsubsection{LIME (Local Interpretable Model-agnostic Explanations)}
LIME creates a surrogate model $g$ (e.g., a linear regression) that is locally faithful to the complex model $f$ around the instance $x$ \cite{ribeiro2016why}. It generates a neighborhood of perturbed samples $\mathcal{Z}$ weighted by a proximity kernel $\pi_x$. The explanation is found by minimizing the loss $\mathcal{L}$:
\begin{equation}
    \xi(x) = \operatorname*{argmin}_{g \in G} \mathcal{L}(f, g, \pi_x) + \Omega(g)
\end{equation}
where $\Omega(g)$ penalizes the complexity of the explanation. While LIME is versatile, its reliance on random sampling introduces instability, where repeated runs on the same input can yield different explanations \cite{hryniewska2024normensemblexai}.

\subsubsection{SHAP (SHapley Additive exPlanations)}
SHAP adapts cooperative game theory to XAI, treating features as players in a coalition. The attribution of a feature is its Shapley value - the average marginal contribution of the feature value across all possible feature coalitions \cite{lundberg2017unified}. The SHAP value $\phi_i$ is defined as:
\begin{equation}
    \phi_i(f, x) = \sum_{z' \subseteq x' \setminus \{i\}} \frac{|z'|! \, (M - |z'| - 1)!}{M!} \bigl[f_x(z' \cup \{i\}) - f_x(z')\bigr]
\end{equation}
SHAP values are the unique additive attribution satisfying \textit{Local Accuracy}, \textit{Missingness}, and \textit{Consistency} \cite{lundberg2017unified}, thereby inheriting the classical Shapley axioms of Efficiency, Symmetry, Dummy, and Additivity. However, exact computation is exponential in the number of features, leading to approximation methods like GradientSHAP or KernelSHAP, which reintroduce estimation noise.

\subsection{Localization Metrics}
\label{subsec:localization-metrics}

Localization metrics assess whether the regions highlighted by an attribution map correspond to semantically meaningful areas of the input. Unlike Faithfulness metrics, which probe the model's response to perturbation, Localization metrics require an external ground-truth mask delineating the object or region of interest. Two widely adopted metrics in this family are the Pointing Game and Relevance Mass Accuracy.

\subsubsection{Pointing Game}
The Pointing Game evaluates whether the single most salient pixel in the attribution map falls within the annotated region of interest. Given an explanation $E$ and a binary segmentation mask $\mathcal{M}$, the Pointing Game score for a single image is defined as:
\begin{equation}
    \text{PG} = \begin{cases} 1 & \text{if } \operatorname{argmax}_{i}|E_i| \in \mathcal{M}, \\ 0 & \text{otherwise.} \end{cases}
\end{equation}
The metric is typically averaged across the test set to produce a hit rate. Its binary nature makes it interpretable but coarse: a method that assigns high attribution to many pixels inside the mask but whose global maximum lies just outside receives a score of zero. The Pointing Game in its present form was introduced by \citet{pg} as part of the Excitation Backprop study, building on the class-activation-mapping line of work that first established the use of localization maps for weakly supervised tasks.

\subsubsection{Relevance Mass Accuracy}
Relevance Mass Accuracy (RMA) captures the fraction of total positive attribution mass that falls within the segmentation mask:
\begin{equation}
    \text{RMA} = \frac{\sum_{i \in \mathcal{M}} \max(E_i, 0)}{\sum_{i} \max(E_i, 0)}.
\end{equation}
RMA ranges from $0$ (no positive attribution inside the mask) to $1$ (all positive attribution is localized). Unlike the Pointing Game, RMA evaluates the entire attribution distribution rather than a single pixel, providing a smoother and more discriminative signal \cite{hedstrom2023quantus}.

\subsubsection{Role of Segmentation Masks}
Both Pointing Game and RMA require a binary segmentation mask aligned with the input image. This requirement limits Localization evaluation to datasets that provide pixel-level annotations. The ISIC 2017 dataset (Task~1) includes expert-annotated lesion boundaries for every image, enabling Localization assessment without any additional labeling effort. Many commonly used benchmarks lack such masks, which is one reason Localization metrics are underreported in the literature despite their intuitive appeal.

\section{The Crisis of Evaluation}
\label{sec:crisis}

A critical challenge in XAI research is the absence of "ground truth." Unlike classification tasks where a label $y$ is known, there is rarely a verified map of which pixels a neural network \textit{actually} used. This ambiguity led to a reliance on "Plausibility" - visual assessment by humans - which proved deceptive.

A seminal study by Adebayo et al., titled "Sanity Checks for Saliency Maps," demonstrated that visual quality does not imply faithfulness \cite{adebayo2018sanity}. They introduced the \textbf{Model Parameter Randomization Test}, where model weights are progressively randomized from the output layer to the input layer. A valid explanation method should degrade to noise as the model's knowledge is destroyed.
\begin{itemize}
    \item \textbf{Failure of Guided Backpropagation:} The study revealed that Guided Backpropagation and Guided Grad-CAM produced nearly identical heatmaps for a trained model and a random network. This indicates these methods act largely as edge detectors, independent of the model's learned representations.
    \item \textbf{Implication:} This finding established a "Crisis of Evaluation" \cite{hedstrom2023quantus}, necessitating a shift from human-grounded evaluation (which measures trust/plausibility) to functionally-grounded evaluation (which measures mathematical properties like fidelity).
\end{itemize}

\section{Taxonomy of Evaluation Metrics}
\label{sec:taxonomy}

To rigorously assess FAE models, recent literature proposes a modular taxonomy of metrics. This thesis adopts the framework formalized by Quantus \cite{hedstrom2023quantus}, covering all six of its categories: Faithfulness, Robustness, Localization (Section~\ref{subsec:localization-metrics}), Complexity, Randomization, and Axiomatic correctness. Figure~\ref{fig:metric-taxonomy} maps the twelve concrete metrics evaluated in this thesis onto these categories; each is defined in the corresponding subsection below, and their configuration is fixed in Chapter~\ref{ch:methodology}.

\begin{figure}[htbp]
    \centering
    \begin{tikzpicture}[
        card/.style={draw, rounded corners, align=left, inner sep=5pt,
                     text width=0.28\linewidth, font=\footnotesize},
        root/.style={draw, rounded corners, align=center, inner sep=5pt,
                     text width=0.55\linewidth, font=\small\bfseries,
                     fill=black!5},
        arr/.style={-{Latex[length=1.6mm]}, semithick}
    ]
        \node[root] (root) {XAI evaluation metrics (Quantus taxonomy)};

        \node[card, below left=7mm and 0.32\linewidth of root.south,
              anchor=north] (faith)
            {\textbf{Faithfulness}\\[1pt]
             Faithfulness Corr.\ ($\uparrow$)\\Pixel Flipping ($\downarrow$)};
        \node[card, below=7mm of root.south, anchor=north] (rob)
            {\textbf{Robustness}\\[1pt]
             Max-Sensitivity ($\downarrow$)\\Avg-Sensitivity ($\downarrow$)};
        \node[card, below right=7mm and 0.32\linewidth of root.south,
              anchor=north] (loc)
            {\textbf{Localization}\\[1pt]
             Pointing Game ($\uparrow$)\\Relevance Mass Acc.\ ($\uparrow$)};
        \node[card, below=4mm of faith] (comp)
            {\textbf{Complexity}\\[1pt]
             Sparseness ($\uparrow$)\\Complexity ($\downarrow$)};
        \node[card, below=4mm of rob] (rand)
            {\textbf{Randomization}\\[1pt]
             Random Logit ($\downarrow$)\\MPRT ($\downarrow$)};
        \node[card, below=4mm of loc] (axio)
            {\textbf{Axiomatic}\\[1pt]
             Completeness\\Non-Sensitivity ($\downarrow$)};

        \foreach \c in {faith, rob, loc}
            \draw[arr] (root.south) -- (\c.north);
    \end{tikzpicture}
    \caption{The six Quantus metric categories \cite{hedstrom2023quantus} and
    the twelve concrete metrics evaluated in this thesis, two per category.
    Arrows in parentheses give each metric's score direction
    ($\uparrow$~higher is better, $\downarrow$~lower is better);
    Completeness is a binary axiom check. The redundancy analysis of
    Chapter~\ref{ch:experiments} shows that not all twelve carry independent
    signal.}
    \label{fig:metric-taxonomy}
\end{figure}

\subsection{Faithfulness (Fidelity)}
Faithfulness quantifies how accurately the explanation reflects the model's internal decision logic, regardless of human intelligibility.

\subsubsection{Pixel Flipping (AUC)}
Pixel Flipping, introduced by \citet{bach2015lrp} and generalized to region
perturbation with an area-under-the-curve protocol by \citet{pixelflip},
evaluates faithfulness by iteratively removing the most important features (as
ranked by the attribution map) and measuring the resulting degradation in the
model's output.
Pixels are replaced in descending order of attribution magnitude with a baseline value (e.g., black), producing a curve of model outputs as ever more of the highest-ranked features are removed. The metric summarizes this degradation curve by its area under the curve:
\begin{equation}
    \text{PixelFlipping}(x) = \frac{1}{N} \sum_{k=1}^{N} f\bigl(x^{(k)}\bigr),
\end{equation}
where $x^{(k)}$ is the image with the top $k$ most important features replaced. A faithful attribution identifies the truly critical features first, so the prediction collapses early and the area under the degradation curve is small: \emph{lower} scores are better. The Quantus implementation groups features into steps of configurable size (\texttt{features\_in\_step}), making the metric tractable for high-resolution inputs \cite{hedstrom2023quantus}.

This formulation generalizes Deletion and Insertion curves \cite{petsiuk2018rise}: Deletion corresponds to removal in descending attribution order, while Insertion reverses the process by adding features to a blank canvas.

A structural caveat applies to this whole family of perturbation-based metrics: replacing pixels with a baseline produces inputs that lie off the training distribution, so part of the measured prediction drop may reflect the model's reaction to out-of-distribution inputs rather than the quality of the attribution. Recent work addresses this confound explicitly---the F-Fidelity framework removes it via fine-tuning and random masking, recovering ground-truth explainer rankings that the naive protocol distorts \citep{zheng2025ffidelity}.

\subsubsection{Faithfulness Correlation}
Faithfulness Correlation, as implemented in Quantus, measures the Pearson correlation between attribution magnitude and the change in model output when features are masked in random subsets \cite{hedstrom2023quantus}. Given \texttt{nr\_runs} randomly sampled subsets $S_1, \dots, S_R$ of features:
\begin{equation}
    \text{FC} = r\!\left(\left\{\sum_{i \in S_j} E_i\right\}_{j=1}^{R},\; \left\{f(x) - f(x_{\setminus S_j})\right\}_{j=1}^{R}\right),
\end{equation}
where $r(\cdot, \cdot)$ denotes the Pearson correlation coefficient and $x_{\setminus S_j}$ is the input with features in $S_j$ replaced by a baseline. A high positive correlation ($\text{FC} \to +1$) indicates that regions with large attribution values are indeed the regions whose removal most degrades the model's prediction.

A related measure, Infidelity \citep{yeh2019}, computes the expected squared error $\mathbb{E}[(I^T E - \Delta f)^2]$ between the attribution-weighted perturbation and the actual output change. Although theoretically motivated, Infidelity is not used in the evaluation framework of this thesis; Faithfulness Correlation provides a comparable signal while being more interpretable and directly comparable across methods.

\subsection{Robustness (Stability)}
\label{subsec:robustness}
Robustness ensures that slight, non-semantic changes to the input (e.g., Gaussian noise) do not cause disproportionate changes in the explanation. This is crucial for user trust and defense against adversarial attacks.

The theoretical foundation for explanation robustness is Lipschitz continuity: a method $\Phi$ is robust if the ratio $\|\Phi(x) - \Phi(x')\| / \|x - x'\|$ remains bounded for inputs $x'$ in a neighborhood of $x$ \cite{alvarez2018robustness}. The metrics adopted in this thesis operationalize this principle via sensitivity analysis under random perturbations.

\subsubsection{Max-Sensitivity}
Max-Sensitivity, proposed by \citet{yeh2019} in ``On the (In)fidelity and Sensitivity of Explanations,'' measures the worst-case explanation change under small random perturbations of the input,
$\mathrm{SENS}_{\mathrm{MAX}} = \max_{\|\delta\| \le r} \|\Phi(x + \delta) - \Phi(x)\|$.
The Quantus implementation used in this thesis estimates it from $N$
perturbation samples $\delta_1, \dots, \delta_N$ drawn uniformly with
magnitude bounded by \texttt{lower\_bound}, and reports the change relative to
the unperturbed attribution:
\begin{equation}
    \text{MaxSens}(\Phi, x) = \max_{j = 1, \dots, N} \frac{\|\Phi(x + \delta_j) - \Phi(x)\|_F}{\|\Phi(x)\|_F}.
\end{equation}
The implementation uses \texttt{nr\_samples}$= 10$ perturbation draws with \texttt{lower\_bound}$= 0.2$ controlling the perturbation magnitude. Lower values indicate a more stable explanation method.

\subsubsection{Avg-Sensitivity}
Avg-Sensitivity replaces the maximum with the mean across the same set of perturbations:
\begin{equation}
    \text{AvgSens}(\Phi, x) = \frac{1}{N} \sum_{j=1}^{N} \frac{\|\Phi(x + \delta_j) - \Phi(x)\|_F}{\|\Phi(x)\|_F}.
\end{equation}
Lower is better. Together, Max-Sensitivity and Avg-Sensitivity provide complementary views: the former captures the tail behavior (susceptibility to worst-case perturbations), while the latter captures the expected behavior. Both are implemented in Quantus \cite{hedstrom2023quantus} and share hyperparameters in this thesis for comparability.

\subsection{Complexity and Sparsity}
\label{subsec:complexity}
Complexity metrics adhere to the principle of Occam's Razor: simpler explanations are generally preferred for human cognition.
\begin{itemize}
    \item \textbf{Sparseness:} Quantifies the concentration of the attribution distribution using a Gini-like coefficient. A perfectly sparse explanation assigns all relevance to a single feature (score $\to 1$), while a uniform explanation scores $0$. Higher is better.
    \item \textbf{Complexity:} Calculated as the Shannon entropy of the normalized absolute attribution distribution. High entropy indicates a diffuse explanation, while low entropy indicates a focused, sparse explanation. Lower is better.
\end{itemize}
Both metrics are computed directly from the attribution map without model access, making them applicable to any FAE method regardless of its generation mechanism \cite{hedstrom2023quantus}.

\subsection{Randomization}
\label{subsec:randomization}
Randomization tests assess whether an explanation method genuinely relies on learned model parameters rather than spurious input statistics. If an explanation remains unchanged after the model's parameters are destroyed, the method is not faithfully reflecting the model's decision process.

\subsubsection{Model Parameter Randomisation Test (MPRT)}
Proposed by Adebayo et al.\ as part of the ``Sanity Checks for Saliency Maps'' study \cite{adebayo2018sanity}, MPRT progressively randomizes the model's weights layer by layer, proceeding from the output layer toward the input (top-down order). After each randomization step, the explanation is recomputed and its similarity to the original explanation is measured via Spearman or Pearson correlation. The metric reports the average correlation across all randomization depths:
\begin{equation}
    \text{MPRT}(\Phi) = \frac{1}{L} \sum_{\ell=1}^{L} \rho\bigl(\Phi(f, x),\; \Phi(f^{(\ell)}, x)\bigr),
\end{equation}
where $f^{(\ell)}$ denotes the model with the top $\ell$ layers randomized and $\rho$ is the chosen correlation measure. Lower average correlation indicates a better method---one that is genuinely sensitive to the model's learned parameters. As discussed in Section~\ref{sec:crisis}, Guided Backpropagation and Guided Grad-CAM fail this test, producing nearly identical maps for trained and randomized models.

\subsubsection{Random Logit Test}
The Random Logit test complements MPRT by perturbing the explanation target rather than the model weights. The target class index is replaced with a randomly sampled alternative class, and the resulting explanation is compared to the original. The metric computes the correlation between the two explanation maps:
\begin{equation}
    \text{RandomLogit}(\Phi) = \rho\bigl(\Phi(f, x, c),\; \Phi(f, x, c')\bigr), \quad c' \neq c.
\end{equation}
A lower correlation indicates that the explanation is genuinely class-discriminative---it changes when the target class changes---rather than merely reflecting generic input features. Both MPRT and Random Logit are implemented in Quantus \cite{hedstrom2023quantus}.

\subsection{Axiomatic}
\label{subsec:axiomatic}
Axiomatic metrics verify whether explanations satisfy formal mathematical axioms that a ``correct'' attribution should obey. These axioms provide theoretical guarantees that hold independently of the specific model or input.

\subsubsection{Completeness}
The Completeness axiom, formalized by Sundararajan et al.\ \cite{sundararajan2017axiomatic}, requires that the sum of individual feature attributions equals the difference between the model output at the input $x$ and a baseline $x'$:
\begin{equation}
    \left|\sum_{i} E_i - \bigl(f(x) - f(x')\bigr)\right| \approx 0.
\end{equation}
The Quantus implementation used here does not score the magnitude of the deviation: it returns a per-sample Boolean, $1$ where the summed attribution equals $f(x) - f(x')$ under strict floating-point equality and $0$ otherwise (higher is better). As Section~\ref{sec:exp-redundancy} reports, that equality never holds in practice, so the metric is constant at $0$ and carries no ranking signal. Completeness is only meaningful for methods that explicitly satisfy or approximate the completeness property by construction---specifically Integrated Gradients, DeepLift, and LRP. For other methods, the metric returns NaN and is excluded from aggregation.

\subsubsection{Non-Sensitivity}
Non-Sensitivity measures how much attribution a method assigns to features that are provably non-contributory (i.e., features whose removal does not change the model output). A faithful attribution should assign zero relevance to such features. The metric evaluates this by performing a forward pass for each feature (or a downsampled subset thereof) to identify non-contributory features, then measuring the attribution mass on those features. Lower is better.

Non-Sensitivity is computationally intensive, as it requires $O(d)$ forward passes, where $d$ is the input dimensionality. In practice, the input is downsampled (e.g., to $56 \times 56$) to make the computation tractable \cite{hedstrom2023quantus}. Due to this cost, Non-Sensitivity is typically excluded from routine evaluations and reserved for targeted diagnostics.

\section{Ensemble XAI: Theoretical Foundations}

Given the trade-offs between these metrics - where Gradient methods often excel in Faithfulness but fail in Robustness, and LIME excels in Local Interpretability but fails in Stability - single-method reliance is increasingly viewed as insufficient.

Ensembling in XAI draws inspiration from predictive ensembles (e.g., Random Forests), operating on the hypothesis that aggregating uncorrelated estimators reduces variance and bias \cite{hryniewska2024normensemblexai}.
\begin{equation}
    E_{\text{ensemble}} = \mathcal{A}(\mathcal{N}(\Phi_1(x)), \mathcal{N}(\Phi_2(x)), \dots, \mathcal{N}(\Phi_n(x)))
\end{equation}
where $\mathcal{N}$ is a normalization function and $\mathcal{A}$ is an aggregation function.

Hryniewska-Guzik et al. identify that the primary challenge in XAI ensembling is the heterogeneous scale of attributions (e.g., SHAP values are additive log-odds, while Grad-CAM is feature activation) \cite{hryniewska2024normensemblexai}. They propose \textbf{NormEnsembleXAI}, which utilizes robust normalization techniques such as Second Moment Scaling to align these diverse explanations before aggregation via Mean, Max, or Minimum functions. Subsequent work strengthens the ensembling premise: \citet{decker2024provably} show that an \emph{optimized} convex combination of attribution methods provably improves explanation quality over any single constituent, replacing NormEnsembleXAI's fixed combinators with learned weights. This thesis builds upon the NormEnsembleXAI framework to evaluate whether ensembling can simultaneously optimize the conflicting objectives of Faithfulness and Robustness.

It is important to distinguish two levels of aggregation that appear in this thesis. \emph{Attribution ensembling}, as performed by NormEnsembleXAI, combines multiple attribution maps $\Phi_1(x), \dots, \Phi_n(x)$ into a single consolidated explanation $E_{\text{ensemble}}$ (a third, orthogonal axis---aggregating explanations across \emph{models} rather than across methods---has also been explored \citep{vascotto2025multimodel}). \emph{Metric aggregation}, which constitutes this thesis's novel contribution, operates at a different level: it combines multiple evaluation metric scores into a single effectiveness index $E(\Phi)$ that characterizes how well a given FAE method performs across all quality dimensions. Chapter~\ref{ch:methodology} develops the metric-aggregation layer on top of the NormEnsembleXAI normalization framework, treating both individual and ensembled attributions as inputs to the same multi-criteria evaluation pipeline.

\section{Multi-Metric Aggregation and Evaluation Pipelines}
\label{sec:multi-metric-aggregation}

The taxonomy presented in Section~\ref{sec:taxonomy} provides a principled
vocabulary for assessing individual properties of explanations. In practice,
however, no single metric suffices to characterize an FAE method: Faithfulness
captures fidelity to the model, Robustness captures stability under
perturbation, and Complexity captures cognitive tractability---and these axes
are neither independent nor unidirectional. A method may score favorably on one
while deteriorating on another, producing rankings that are sensitive to the
choice of metric \citep{krishna2024disagreement, hedstrom2024metaquantus}. This
section addresses three questions that arise once multiple metrics enter the
evaluation: (i)~whether the available metrics are mutually informative or
merely redundant, (ii)~how to validate the metrics themselves, and (iii)~how
to aggregate non-redundant, validated metrics into a coherent scalar
assessment suitable for method comparison.

\subsection{The Metric Redundancy Problem}
\label{subsec:redundancy}

The proliferation of evaluation metrics within libraries such as Quantus
\citep{hedstrom2023quantus} raises a concern that has attracted growing
attention: the superset of available metrics contains substantial redundancy,
which inflates the apparent rigor of an evaluation without adding information.
Formally, given a set of $K$ metrics $\{M_1, M_2, \dots, M_K\}$ applied to an
explanation $\Phi(f, x)$, one can compute the empirical correlation matrix
across a test set $\mathcal{D}$:
$$
\rho_{ij} \;=\; \operatorname{corr}_{x \sim \mathcal{D}}\!\Bigl(
    M_i\bigl(\Phi, f, x\bigr),\; M_j\bigl(\Phi, f, x\bigr)
\Bigr).
$$
\citet{tomsett2020sanity} were among the first to quantify this empirically,
reporting that metrics within the same category---Faithfulness in
particular---exhibit surprisingly low pairwise correlation when applied to
image classifiers, despite purporting to measure the same underlying property.
This finding has since been replicated across modalities:
\citet{bommer2024finding} observed metric-dependent rankings of XAI methods in
a climate-science benchmark, and \citet{klein2024navigating} demonstrated
systematic rank reversals between faithfulness operationalizations on a
large-scale computer vision benchmark. At the level of individual metrics,
\citet{barr2023disagreement} document an analogous disagreement problem among
faithfulness metrics themselves---the metric-level counterpart of the
explainer-level disagreement of \citet{krishna2024disagreement}. \citet{pawlicki2024necessity} argue that
this proliferation is counterproductive: many widely-used metrics either
duplicate information already captured by simpler alternatives or are
ineffective by their own operational definition, and recommend a deliberate
reduction to a representative subset.

Two practical consequences follow. First, reporting scores across all metrics
in a category overstates evidence: a method that wins on five correlated
metrics has, in effect, won once. Second, aggregation schemes that treat
metrics as independent voters (e.g., averaging ranks across all available
faithfulness metrics) are biased toward whichever operational family is most
heavily represented in the chosen library. A principled evaluation must
therefore begin with a redundancy analysis and proceed with a reduced,
justified subset. Meta-reviews of the metric landscape reach the same
conclusion from the reporting side: \citet{gipiskis2026cards} find that
correlations between related metrics are rarely reported at all, and propose
evaluation cards that make them a required reporting field.

\subsection{Meta-Evaluation: Validating the Validators}
\label{subsec:meta-evaluation}

The redundancy problem reveals that metrics disagree with each other; a deeper
problem is that it is not always clear which metric is \emph{correct}. Since
no ground-truth explanation exists against which to calibrate a metric, the
standard evaluation paradigm lacks the reference signal used elsewhere in
supervised machine learning. \citet{hedstrom2024metaquantus} formalize this as
the \emph{meta-evaluation problem}: how does one evaluate a metric that is
itself supposed to evaluate an explanation? Their proposed framework,
MetaQuantus, subjects each metric to two stress tests grounded in plausible
failure modes:

\begin{itemize}
    \item \textbf{Noise Resilience (NR)}: under controlled, semantically
    meaningless perturbations of model inputs or weights, a reliable metric
    should yield stable scores. High score variance under noise indicates
    that the metric is responding to spurious signal.
    \item \textbf{Adversarial Reactivity (AR)}: under controlled degradation
    of the explanation itself (e.g.,\ partial randomization), a reliable
    metric should show a monotonic drop in score. A metric that remains
    unchanged under explanation degradation is not tracking explanation
    quality.
\end{itemize}

A metric that performs poorly on either test is deprioritized as an unreliable
signal, regardless of its theoretical motivation. (This thesis implements
MetaQuantus-\emph{inspired} estimators of these two properties rather than
the published estimators; the operationalization is specified in
Chapter~\ref{ch:methodology}.) The concern is not merely
theoretical: evaluating metrics against transparent models whose ground-truth
explanations are known, \citet{miro2025fidelity} conclude that current
fidelity metrics are not yet reliable enough for real-world scenarios, which
underscores the need to validate a metric before trusting its verdict, and
\citet{fresz2024classification} likewise treat the reliability of the
evaluation metrics themselves---their objectivity and consistency in the
psychometric sense---as a first-class requirement for trustworthy XAI
evaluation. This
meta-evaluation layer
is essential for the framework proposed in this thesis, as it provides a
principled criterion for the \emph{weighting} of metrics in the aggregation
stage of Chapter~\ref{ch:methodology},
complementing the correlation-based redundancy analysis of
Section~\ref{subsec:redundancy}.

\subsection{Normalization and Aggregation Strategies}
\label{subsec:aggregation-strategies}

Once a non-redundant, meta-validated subset $\{M_1, \dots, M_K\}$ is selected,
the individual scores must be combined into a scalar that supports method
comparison. Two technical obstacles arise: heterogeneous scales and
heterogeneous directions. A Faithfulness Correlation score lies in $[-1, 1]$;
a Max-Sensitivity score is unbounded above and lower-is-better; a Sparseness
score is in $[0, 1]$ and higher-is-better. Direct averaging is meaningless
without rectification.

Let $M_k(\Phi)$ denote the raw score of metric $k$ for explanation $\Phi$, and
let $s_k \in \{+1, -1\}$ encode the direction convention ($+1$ for
higher-is-better). A normalized score $\tilde{M}_k$ is obtained via
$$
\tilde{M}_k(\Phi) \;=\; s_k \cdot \psi_k\!\bigl(M_k(\Phi)\bigr),
$$
where $\psi_k$ is a scale-rectifying transformation. Three choices dominate
the literature:
\begin{enumerate}
    \item \textbf{Min--max normalization}:
    $\psi_k(x) = (x - \min_k) / (\max_k - \min_k)$, mapping scores to $[0, 1]$.
    Sensitive to outliers.
    \item \textbf{$z$-score normalization}:
    $\psi_k(x) = (x - \mu_k) / \sigma_k$, centering on the mean with unit
    variance.
    \item \textbf{Second Moment Scaling}:
    $\psi_k(x) = x / \sqrt{\mathbb{E}[x^2]}$, employed by
    \citet{hryniewska2024normensemblexai} in the NormEnsembleXAI framework to
    align explanations of heterogeneous origin (e.g.,\ SHAP's log-odds
    contributions vs.\ Grad-CAM's activation magnitudes) prior to ensemble
    aggregation.
\end{enumerate}

Once normalized, aggregation proceeds via a scalar combinator $\mathcal{A}$,
yielding the effectiveness index
$$
E(\Phi) \;=\; \mathcal{A}\!\Bigl(
    \tilde{M}_1(\Phi),\; \tilde{M}_2(\Phi),\; \dots,\; \tilde{M}_K(\Phi)
\Bigr).
$$
Common choices for $\mathcal{A}$ include the weighted arithmetic mean
$\sum_k w_k \tilde{M}_k$, the minimum $\min_k \tilde{M}_k$ (a conservative
operator favoring methods acceptable on all criteria), and the geometric mean
(penalizing methods that fail severely on any single axis). The minimum
operator corresponds to evaluating a method by its worst dimension and has a
natural interpretation in safety-critical domains: a method that is faithful
but unstable is ineligible regardless of how high its faithfulness score is.
Beyond such fixed combinators, recent work frames the combination step as a
multi-criteria decision-making (MCDM) problem: \citet{chatterjee2025mcdm}
weight competing quality metrics with TOPSIS-style operators and aggregate
explainer \emph{rankings} rather than scores, offering a rank-based
alternative to the score-based weighted mean adopted in this thesis.

\subsection{Evaluation Pipelines and the Autoweighted System}
\label{subsec:evaluation-pipelines}

The framework introduced above---(i)~redundancy-aware metric selection,
(ii)~meta-validation, (iii)~direction rectification and scale normalization,
(iv)~aggregation---constitutes an \emph{evaluation pipeline}. The pipeline
abstraction is significant because it isolates configuration decisions (which
metrics, which normalization, which aggregator, which weights) from the XAI
method under evaluation, enabling reproducible, principled comparison across
studies.

A key design choice within this pipeline is the weight vector
$\mathbf{w} = (w_1, \dots, w_K)$. Fixed equal weighting assumes all retained
criteria contribute equally---a strong assumption that the redundancy analysis
alone does not validate. \citet{hryniewska2024normensemblexai} introduce the
\emph{Autoweighted} mechanism, a data-adaptive weighting in which each
component's weight is derived from the statistics of its own scores rather
than set by hand, so that dominant, trustworthy signals receive higher weight
without manual tuning. This thesis adapts the Autoweighted idea to the
metric-aggregation layer by weighting each metric by the inverse coefficient
of variation of its scores across the test set---metrics whose scores are
internally consistent are trusted more than metrics whose scores are
noisy---with the exact formulation given in Chapter~\ref{ch:methodology}
(Equation~\eqref{eq:autoweighted}). Crucially, the thesis keeps this
Autoweighted scheme intact and \emph{separately} formulates a
MetaQuantus-discounted extension (Equation~\eqref{eq:mq-discount}), in which
each metric's Autoweighted coefficient is attenuated by its meta-evaluation
reliability (Noise Resilience and Adversarial Reactivity from
\citet{hedstrom2024metaquantus}). Keeping the two schemes distinct is what
allows the experiments of Chapter~\ref{ch:experiments} to isolate the effect
of the reliability discount on the resulting FAE ranking.

Recent work has emphasized that evaluation pipelines are themselves a
compliance instrument. \citet{seth2025bridging} argue that standardized,
reproducible pipelines---rather than ad-hoc metric selection---are a
prerequisite for operationalizing the transparency requirements of the EU AI
Act, aligning the technical contribution of this thesis with the regulatory
motivation established in Section~\ref{sec:motivation}.